\subsection{正埃尔米特型}

定义 2. --- 设 E 为域 K 上的向量空间。E 上的埃尔米特型 f 称为正的，如果 \(f(x, x) \geqslant 0\) 对所有 \(x \in E\) 成立。

显然，向量空间 E 上的全体埃尔米特型构成域 R 上的一个向量空间（当 K 是 C 时，不构成域 C 上的向量空间）：在这个空间中，由定义 2 与命题 1，诸正埃尔米特型构成一个真尖凸锥（II, p.~10）。

命题 2. --- 若 \(f\) 是正埃尔米特型，则有

\[
\left| f (x, y) \right| ^ {2} \leqslant f (x, x) f (y, y)\tag{7}
\]

对 E 中每个 x 与 y（Cauchy-Schwarz 不等式）。

先设 \(f(y, y) \neq 0\) 。对每个 \(\xi \in \mathbf{K}\) ，有

\[
f (y, y) f (x + \xi y, x + \xi y) \geqslant 0
\]

它可写成

\[
f (x, x) f (y, y) - | f (x, y) | ^ {2} + (\xi f (y, y) + \overline {{f (x , y)}}) (\bar {\xi} f (y, y) + f (x, y)) \geqslant 0.
\]

在这个不等式中把 \(\xi\) 换成 \(-\overline{f(x,y)} / f(y,y)\) ，即得 (7)。若 \(f(x,x)\neq 0\) ，论证类似。

最后，若 \(f(x, x) = f(y, y) = 0\) ，则 \(f(x + \xi y, x + \xi y) \geqslant 0\) 对所有 \(\xi \in \mathbf{K}\) 成立，这可写成

\[
\xi f (x, y) + \overline {{{{\xi f (x , y)}}}} \geqslant 0.
\]

在这个不等式中把 \(\xi\) 换成 \(-\overline{f(x,y)}\) ，得 \(-2|f(x,y)|^2 \geqslant 0\) ，从而 \(f(x,y) = 0\) ；在此情形我们同样得到 (7)。

推论 1. --- 若 \(f\) 是正埃尔米特型，则集 \(\mathbf{N}\) —— 即全体 \(x \in \mathbf{E}\) ，使 \(f(x, x) = 0\) —— 与全体 \(x \in \mathbf{E}\) （使 \(f(x, y) = 0\) 对所有 \(y \in \mathbf{E}\) 成立的）组成的向量子空间重合。

推论 2. --- 为使一个正埃尔米特型非退化，必要且充分的是关系 \(x \neq 0\) 蕴含 \(f(x, x) > 0\) 。

这由推论 1 立即得出。

对 E 上的每个正埃尔米特型 f，与 \(f(\mathrm{V}, \mathrm{p.~2})\) 相伴的非退化埃尔米特型显然是 E/N 上的正埃尔米特型。

命题 3. --- 设 f 为 E 上的正埃尔米特型。令

\[
p (x) = f (x, x) ^ {1 / 2}
\]

对所有 \(x \in \mathbb{E}\) 。则 \(p\) 是 \(\mathbb{E}\) 上的半范数，且是范数当且仅当 \(f\) 非退化。只需证明不等式 \(p(x + y) \leqslant p(x) + p(y)\) 。但我们有

\[
f (x + y, x + y) = f (x, x) + f (y, y) + f (x, y) + \overline {{f (x , y)}}
\]

且由 Cauchy-Schwarz 不等式

\[
\begin{array}{c} f (x + y, x + y) \leqslant f (x, x) + f (y, y) + 2 (f (x, x) f (y, y)) ^ {1 / 2} \\ = \bigl (f (x, x) ^ {1 / 2} + f (y, y) ^ {1 / 2} \bigr) ^ {2}. \end{array}\tag{8}
\]

注. --- 1) 假设 \(f\) 是正的且非退化的，并设 \(x, y\) 是两个 \(\neq 0\) 的向量。Cauchy-Schwarz 不等式的证明表明，若 (7) 的两端相等，则存在标量 \(\xi\) 使 \(f(x + \xi y, x + \xi y) = 0\) ，从而 \(x + \xi y = 0\) ，换言之，x 与 y 线性相关；反之是直接的。不等式 (8) 的证明表明，等式 \(p(x + y) = p(x) + p(y)\) 只当 x 与 y 线性相关时才可能成立；若 \(y = \lambda x\) ，前述等式可写成 \(|1 + \lambda| = 1 + |\lambda|\) ，它蕴含 \(\lambda\) 是实数且是正的。

\begin{enumerate}
\def\labelenumi{\arabic{enumi})}
\setcounter{enumi}{1}
\tightlist
\item
  设 \(f\) 为 \(\mathbf{E}\) 上的正埃尔米特型，并给 \(\mathbf{E}\) 赋以半范数 \(x \mapsto f(x, x)^{1/2}\) ；若 \(\dot{f}\) 是定义在 \(\mathrm{E/N}\) 上的与 \(f\) 相伴的正的非退化埃尔米特型，则把范数 \(x \mapsto \hat{f}(\dot{x}, \dot{x})^{1/2}\) 赋予 \(\mathrm{E/N}\) 所得的赋范空间就是与 \(\mathbf{E}\) 相伴的赋范空间（II, p.~5）。
\end{enumerate}

定义 3. --- 设 E 为域 K 上的向量空间。E 上的半范数 p 称为准希尔伯特的，如果存在 E 上的正埃尔米特型 f，使 \(p(x) = f(x, x)^{1/2}\) 对所有 \(x \in E\) 成立。

注意，对 E 上的半范数 p，至多存在一个正埃尔米特型 f，使 \(p(x) = f(x, x)^{1/2}\) 对所有 \(x \in E\) 成立；这由极化公式得出（V, p.~2）。

